% !PRINTED 532
\weberpage{532}

%<<CONT
eine Gleichung dritten Grades und auf zwei Quadratwurzeln
zurückgeführt werden. Diese Gleichung sechsten Grades ist eine
Totalresolvente, und zwei ihrer Wurzeln genügen zur rationalen
Darstellung der Wurzeln $\alpha$. Die Form dieser Resolventen wird
nicht einfach.

Man kann endlich noch darauf ausgehen, durch Adjunction
einer zur Gruppe $P_3$ gehörigen Function $\xi$ die Function $f(x)$
direct reducibel zu machen und in zwei quadratische Factoren
zu zerlegen. Eine solche Function $\xi$ ist gleichfalls die Wurzel
einer Gleichung sechsten Grades, die sich aber auch durch Imprimitivität
 auf den dritten und zweiten Grad reducirt. Nimmt
man z. B.
\weq{22}{\xi = \alpha + \alpha_2 - \tfrac{1}{2}a_1 = \tfrac{1}{2}\sqrt{v_2},}
so sind von den sechs Wurzeln je zwei entgegengesetzt gleich,
so dass eine cubische Gleichung für $\xi^2$ resultirt. Durch eine
dieser Grössen $\xi$ lassen sich dann die quadratischen Factoren
von $f(x)$ rational darstellen.

Nehmen wir zur Vereinfachung der Formeln $a_1 = 0$ an und
\weq{23}{f(x) = x^4 + ax^2 - bx + c,}
so wird
\[
\begin{aligned}
(\alpha\alpha_2 - \alpha_1\alpha_3)\,\xi &= b \\
\alpha\alpha_2 + \alpha_1\alpha_3 &= \xi^2 + a,
\end{aligned}
\]
also
\weq{24}{2\,\alpha\alpha_2 = \frac{b}{\xi} + \xi^2 + a, \qquad 2\,\alpha_1\alpha_3 = -\frac{b}{\xi} + \xi^2 + a,}
so dass die beiden quadratischen Factoren von $f(x)$ folgende werden:
\weq{25}{
\begin{aligned}
x^2 - \xi x + \tfrac{1}{2}\Bigl(\dfrac{b}{\xi} + \xi^2 + a\Bigr) &= 0 \\
x^2 + \xi x + \tfrac{1}{2}\Bigl(-\dfrac{b}{\xi} + \xi^2 + a\Bigr) &= 0,
\end{aligned}
}
und für $\xi^2$ ergiebt sich aus (24) durch Multiplication die cubische
Gleichung
\[
4c = (\xi^2 + a)^2 - \frac{b^2}{\xi^2}
\]
oder
\weq{26}{\xi^6 + 2a\xi^4 + (a^2 - 4c)\xi^2 - b^2 = 0.}
Nimmt man für $\xi$ irgend eine Wurzel dieser Gleichung
sechsten Grades, so giebt jede der Gleichungen (25) ein Paar
Wurzeln von $f(x) = 0$.
